\needspace{15\baselineskip}% room for the wrapped icon (it cannot break across pages)
\subsection{CO$_2$ battery, 10~h (2,449~\$/kW + 244.9~\$/kWh)}\label{sec:co2-battery}
\techicon{co2-battery}{An inflatable dome holding CO$_2$ gas, beside vessels of liquid CO$_2$ and the compressor-turbine.}
Not part of the default pool of advanced technologies, only listed here for transparency about its exclusion.

Energy Dome's (the company pushing this technology forward) CO$_2$ battery (Figure~\ref{fig:icon-co2-battery}) stores energy by compressing CO$_2$ gas into a liquid and releases it through a turbine.
It is built from conventional compressors, turbines and tanks and suits durations of around 10~hours.
Its landmark is a large white inflatable dome that holds the CO$_2$ as gas at atmospheric pressure while the battery is empty; a 20~MW / 200~MWh plant covers about 5~hectares, most of it dome \citep{utilitydive2024}.

The CAPEX of 2,449~\$/kW plus 244.9~\$/kWh at 10~hours follows Energy Dome's disclosed price of about 220~EUR\textsubscript{2023}/kWh for the whole system, quoted for its first full-scale plant including engineering, construction and margin \citep{esn2023}.
Alliant's US plant of the same size is budgeted at 60--90~million~USD, or 300--450~\$/kWh \citep{utilitydive2024}.

We leave the technology out because it is too similar to iron-air to be included as well.
Both are long-duration storage built from cheap, abundant materials, and both would fill the same gap beyond Li-ion.
Iron-air is the stronger candidate for this role: its energy part costs less than a third (75 against 245~\$/kWh), and it has a large procured project at Pine Island.
The CO$_2$ battery's higher round-trip efficiency (0.75 against 0.41) suits daily cycling better, so it may be assessed as a sensitivity.

The learning rate is 5\,\% on the energy part; there is no cost history.
It is an analogy to mature mechanical storage: the plant consists of conventional compressors, turbines and tanks, whose costs no longer fall much, and pumped hydro, the closest analogue, shows no learning ($-1 \pm 8$\,\%, \citealp{schmidt2018}).

The round-trip efficiency is 0.75 \citep{esn2023}, 0.73--0.75 according to Alliant \citep{utilitydive2024}, and the lifetime 30~years \citep{energydome}.
No fixed O\&M has been published, so the cell is blank.
Existing capacity is the 2.5~MW / 4~MWh demonstrator at Ottana, Sardinia, in trials since 2022 \citep{mittr2022}; the pipeline is a 20~MW / 200~MWh plant at the same site and Alliant Energy's plant of the same size in Columbia, Wisconsin \citep{utilitydive2024}.
The lead time of 2~years is our assumption for a plant of conventional components; the Sardinian plant has slipped from completion in 2024 \citep{utilitydive2024} to 2025 \citep{engie2024}, so it may be longer.
